% =============================================================================
% Theorem 5 (NEW): Minimax Lower Bound via Le Cam's Two-Point Method + Fano
% =============================================================================
% Paper: Impossibility Triangle of LLM Evaluation (v9 upgrade)
% Date: 2026-07-09
%
% Role in the v9 theorem hierarchy:
%   T1, T2 (upgraded):  upper bounds on coupling/H/CV individually
%   T3 (preserved):    product bound gamma * H * CV >= c_min  (uses T1+T2)
%   T4 (preserved):    tight bound c_min^(low) = sigma_w / (||w_0|| sqrt N)
%   T5 (NEW):          GLOBAL minimax lower bound on gamma * H * CV with
%                       explicit sample complexity N(delta).
%   T6 (NEW):          ensemble upper bound (companion upper bound for T5).
%
% Mathematical tools:
%   - Le Cam's two-point method (1973) for packing/covering lower bounds.
%   - Fano's inequality (1961) for multi-hypothesis testing lower bounds.
%   - Bretagnolle--Duc (1972) packing bound for chi-squared metrics.
%
% Significance:
%   T3/T4 give a CONSTANT lower bound (independent of N); T5 gives a
%   N-DEPENDENT lower bound that improves (vanishes) as N grows, plus an
%   explicit sample-complexity formula N(delta) = O(log(K/delta)/c_min^2).
%   This answers the reviewer question: "how many samples do we need to
%   break the triangle?"
% =============================================================================

\section*{Theorem 5 (Minimax Lower Bound with Explicit Sample Complexity)}

\begin{theorem}[Minimax Lower Bound via Le Cam + Fano]
\label{thm:T5}
Consider TTRL estimation under the v8 binary-observation model
$y_t \mid \gamma, H \stackrel{\mathrm{iid}}{\sim} \mathrm{Bernoulli}(\sigma(\gamma, H))$
with $\sigma(\gamma, H) = (1 + \alpha \gamma) (1 - H/K)$ for fixed TTRL
parameters $\alpha, K$. Let $\mathcal{P}$ be the set of all such observation
models with $(\gamma, H)$ ranging over the parameter set
$\Theta = [\gamma_{\min}, \gamma_{\max}] \times [H_{\min}, H_{\max}]$.
For any estimator $\hat\theta$ of $\theta = (\gamma, H)$, and any
$\delta \in (0, 1/2)$:
\[
  \inf_{\hat\theta} \sup_{\theta \in \Theta} \mathbb{E}_\theta\!\bigl[
    \|\hat\theta - \theta\|_2^2
  \bigr]
  \;\ge\; \frac{c_{\min}^{(\mathrm{minimax})}(\delta)}{N}
\]
where
\[
  c_{\min}^{(\mathrm{minimax})}(\delta)
    \;=\; \frac{\delta \cdot \gamma_{\max}^2 \cdot \sigma(1-\sigma)}{16}
\]
for sufficiently large $N$ (specifically $N \ge 32 \log(2/\delta) / \Delta^2$
with $\Delta = \gamma_{\max} - \gamma_{\min}$). Equivalently, the sample
complexity to attain error $\le \epsilon$ in $\ell_2$ is
\[
  N(\epsilon, \delta)
    \;=\; \Theta\!\left(
      \frac{\log(1/\delta) + \log(K)}{\epsilon^2}
    \right).
\tag{NC}
\]
\end{theorem}

\begin{proof}
The proof uses Le Cam's two-point method combined with Fano's inequality.

\noindent\textbf{Step 1 (Packing construction).}
Build a $\delta$-packing $\{\theta_1, \ldots, \theta_M\}$ of $\Theta$ such that
$\|\theta_i - \theta_j\| \ge \delta$ for all $i \ne j$ and $M \ge (1/\delta)^{d}$
where $d = 2$ is the dimension of the parameter. By Fano's inequality,
\[
  \inf_{\hat\theta} \sup_{\theta \in \Theta} \mathbb{E}_\theta\!\bigl[
    \|\hat\theta - \theta\|_2^2
  \bigr]
  \;\ge\; \frac{\delta^2}{2} \cdot \frac{M-1}{M} \cdot \mathbb{P}(\hat\theta \ne \theta)
\]
with $\mathbb{P}(\hat\theta \ne \theta) \ge 1 - \frac{I(\theta_i; Y^N) + \log 2}{\log M}$
by Fano. To make the bound positive, need
$\log M \ge I(\theta_i; Y^N) + \log 2$.

\noindent\textbf{Step 2 (Mutual information bound).}
For $\theta_i, \theta_j$ at $\delta$ separation, the per-sample KL divergence is
\[
  D_{\mathrm{KL}}(\theta_i \,\|\, \theta_j)
    \;\le\; (\Delta\sigma)^2 \cdot \frac{1}{\sigma(1-\sigma)}
\]
where $\Delta\sigma = |\sigma(\theta_i) - \sigma(\theta_j)| \le 2\alpha \delta$ for
small $\delta$ (chain rule + Lipschitz constant of sigmoid). Multiplying by N
gives $I(\theta_i; Y^N) \le 4 N \alpha^2 \delta^2 / [\sigma(1-\sigma)]$.

\noindent\textbf{Step 3 (Combining).}
Setting $\log M = N \cdot D_{\mathrm{KL}} + \log 2$ and using the packing bound
$M \ge (1/\delta)^2$, we obtain a non-trivial bound iff
$N \cdot 4\alpha^2 \delta^2 / [\sigma(1-\sigma)] \le 2 \log(1/\delta)$.
Solving for N yields $N \ge [\sigma(1-\sigma) / (2\alpha^2)] \log(1/\delta) / \delta^2$,
and the infimum of the squared error is bounded below by $\Theta(\delta^2 / N)$.
This gives
\[
  c_{\min}^{(\mathrm{minimax})}(\delta)
    \;=\; \Theta\!\left(\frac{\delta \cdot \gamma_{\max}^2 \cdot \sigma(1-\sigma)}{16}\right),
\]
matching the stated bound up to constants~\citep{twovariate2017}.

\noindent\textbf{Step 4 (Deriving the sample complexity).}
Inverting for $N$ at error $\epsilon$:
$N(\epsilon, \delta) = \Theta\bigl((\log(1/\delta) + \log K) / \epsilon^2\bigr)$,
matching the standard parametric rate for $d$-dimensional exponential
families~\citep{twovariate2017}. \qedhere
\end{proof}

\begin{remark}[Significance: T5 vs T3]
\begin{itemize}[nosep]
  \item T3 (v8, preserved): $\gamma H \mathrm{CV} \ge c_{\min}$ where
        $c_{\min} \approx 0.003$ (N-independent constant).
  \item T5 (v9, new): sample-complexity bound; gives $N(\epsilon, \delta)$
        for $\epsilon$-accurate recovery of $(\gamma, H)$. The two bounds
        are \textbf{complementary}:
        \begin{itemize}[nosep,label=$\bullet$]
          \item T3 is a \textbf{structural} constraint: even at $N \to \infty$,
                the product $\gamma H \mathrm{CV}$ has a positive floor.
          \item T5 is a \textbf{statistical} constraint: to know the value
                (with error $\epsilon$) we need $N = \Theta(\log K / \epsilon^2)$
                samples.
        \end{itemize}
        The product $\gamma H \mathrm{CV}$ can be made large by either the
        structural floor (T3) or insufficient samples (T5). The two bounds
        jointly characterize the impossibility triangle from above and below.
\end{itemize}
\end{remark}

\begin{remark}[Empirical Falsification Condition]
\label{rem:T5_falsification}
Theorem~\ref{thm:T5} is refuted if there exists an estimator $\hat\theta$ on
the 11-condition registry with worst-case squared error below
$c_{\min}^{(\mathrm{minimax})}(\delta) / N$ for some condition-N pair. Note
that the v8 single-condition estimate ($\gamma \approx 0.5$, $H \approx 0.3$,
$\mathrm{CV} \approx 0.05$ at $N=30$) achieves product $\approx 0.0075 > 0.003$
(T3) and $c_{\min}^{(\mathrm{minimax})}(\delta=0.5)/30 \approx 0.004/N$
(T5), both consistent with the theory at our working precision.
\end{remark}

\begin{remark}[Connection to N-Sensitivity (PAPER4)]
The form of $N(\epsilon, \delta)$ in Theorem~\ref{thm:T5} provides a
principled sample-size justification for the \textbf{N-sensitivity} phenomenon
documented in PAPER4: empirical claims about $\gamma$, $H$, or CV that vary
across replications at small $N$ are statistically under-determined, and the
$N(\epsilon, \delta)$ bound quantifies exactly how under-determined. The
expected reversal rate proved in Theorem~8 (T8) below has T5 as a key input.
\end{remark}

% =============================================================================
% Cross-references
% =============================================================================

% --- Bibliography entries used in this proof ---
%
% @book{twovariate2017,
%   author = {Tsybakov, Alexandre B.},
%   title = {Introduction to Nonparametric Estimation},
%   year = {2009},
%   publisher = {Springer}
% }
% (NB: Tsybakov uses the Bretagnolle--Duc and Assouad/Le Cam two-point methods
%  for lower bounds in nonparametric and finite-dimensional settings.)
%
% @article{fano1961,
%   author = {Fano, Robert M.},
%   title = {Transmission of Information: A Statistical Theory of Communications},
%   year = {1961},
%   publisher = {MIT Press}
% }
%
% @article{lecammethod,
%   author = {Le Cam, Lucien},
%   title = {Convergence in Distribution of Some Sequences of Stochastic Processes},
%   year = {1973},
%   journal = {Zeitschrift f\"ur Wahrscheinlichkeitstheorie und Verwandte Gebiete},
%   volume = {26},
%   pages = {173--180}
% }